\documentclass[a4,20pt]{article}

\usepackage{indentfirst}
\usepackage{latexsym}
\usepackage[empty]{fullpage}
\usepackage{titlesec}
\usepackage{marvosym}
\usepackage[usenames,dvipsnames]{color}
\usepackage{verbatim}
\usepackage[hidelinks]{hyperref}
\usepackage[utf8]{inputenc}
\usepackage{fancyhdr}
\usepackage{multicol}
\usepackage{hyperref}
\usepackage{csquotes}
\usepackage{tabularx}
\hypersetup{colorlinks=true,urlcolor=black}
\usepackage[11pt]{moresize}
\usepackage{setspace}
\usepackage{fontspec}
\usepackage[inline]{enumitem}
\usepackage{array}
\newcolumntype{P}[1]{>{\centering\arraybackslash}p{#1}}
\usepackage{anyfontsize}
\usepackage{fontspec}
\usepackage{academicons}
\usepackage{fancyhdr}
\usepackage{fontawesome5}
\usepackage{indentfirst}
\usepackage{lastpage}
\usepackage{parskip}[0.4cm]
\usepackage{hyperref}
\usepackage{url}
%%%% Set Margins
\usepackage[margin=1.75cm, top=1.5cm]{geometry}
\usepackage[english]{babel}
\usepackage[document]{ragged2e}

%%%% Set Fonts
\setmainfont[
BoldFont=EBGaramond-Bold.ttf, %SourceSansPro-Bold.otf
ItalicFont=EBGaramond-Italic.ttf
]{EBGaramond-Regular.ttf}
\setsansfont{EBGaramond-SemiBold.ttf}

%%%% Set Page
\pagestyle{fancy}
\fancyhf{} 
\fancyfoot{}
\renewcommand{\headrulewidth}{0pt}
\renewcommand{\footrulewidth}{0pt}

%%%% Set URL Style
\urlstyle{same}

%%%% Set Indentation
\raggedbottom
\raggedright
\setlength{\tabcolsep}{0in}

%%%% Set Secondary Color
\definecolor{UI_blue}{RGB}{0,0,0}

%%%% Define New Commands


%%%% Set Sections formatting
\titleformat{\section}{
 \color{UI_blue} \scshape \raggedright \large \bfseries
}{}{17.3em}{}[\vspace{-10pt} \hrulefill \vspace{-2pt}]

%%%% Set Subtext Formatting
\newcommand{\subtext}[1]{
#1\par\vspace{-0.2cm}}

\newcommand{\subtextit}[1]{\vspace{0.15cm}
 \textit{ #1 \vspace{-0.2cm}} }
 %%%
\newcommand{\cps}{{\color{red} \quad \textbf{|}\quad}}
\newcommand{\cpshalf}{{\color{red} \hspace{0.2em}\textbf{|}\hspace{0.2em}}}
%%%% Set Item Spacing
\setlist[itemize]{align=parleft,left=10pt..2em}

%%%% New Itemize "Zitemize" Formatting - tighter spacing than itemize
\newenvironment{zitemize}{
\begin{itemize}\itemsep0pt \parskip0pt \parsep1pt}
{\end{itemize}\vspace{-0.5cm}}


%%%% Define Skills Bold Formatting
\newcommand{\hskills}[1]{
\textbf{\bfseries #1} }

%%%% Set Subsection formatting
\titleformat{\subsection}{\vspace{-0.1cm} 
\bfseries \fontsize{11pt}{2cm}}{}{0em}{}[\vspace{-0.2cm}]
%%%%%%

%%%%%%% --------------------------------------------------------------------------------------
\rfoot{{\par\fontfamily{ttf} \selectfont \footnotesize \emph{ \thepage \hspace{1pt} of \pageref{LastPage}}}}
\setlength{\parindent}{4em}
\setlength{\parskip}{0.2cm}
%%%%%%% --------------------------------------------------------------------------------------

\def\Vrulefill#1#2{
    \vspace{-10pt}
    \leavevmode%
    %\hskip-.2in%
    \leaders%
    \vtop{\hsize=.0025in\vskip#1#2}%
    \hfill%
    %\hskip.3in%
    \vspace{-2pt}
}


\begin{document}

%%%%%%% --------------------------------------------------------------------------------------
%%%%%%%  HEADER
%%%%%%% --------------------------------------------------------------------------------------
{{\huge \scshape Statement of Purpose }} %\footnote{Parts of this document were grammatically prepared with the assistance of a generative artificial intelligence tool. The content, however, remains truthful, accurate, and in my own words.}
\hfill { \small \textbf{10 Dec. 2023}} \\
%\hfill {Aeronautics and Astronautics}\\
\Vrulefill{-.05in}{-}
%\rule{10cm}{0.1em}
%%%%%%% --------------------------------------------------------------------------------------
%%%%%%%  Research Interest
%%%%%%% --------------------------------------------------------------------------------------
\justifying

\begin{large} 
Dear Admissions Committee,


With an unwavering commitment to interdisciplinary research, I am driven by a profound passion for leveraging AI solutions in astronomy. I seek to advance astrophysics by developing and employing novel methodologies to tackle the cosmos' most profound mysteries, ranging from the intricate structure of the interstellar medium to the elusive nature of dark matter. I seek to utilize AI as a driving force to unravel celestial mysteries and advance astronomical research. Several labs within Harvard CFA could support me in this research.


As a first-generation student, I have navigated the academic world largely on my own, overcoming challenges and breaking barriers that have fueled my passion for research. My journey into astrophysics and AI has been one of self-discovery and relentless determination, much like an explorer charting a course through the cosmos. Each challenge I’ve overcome and each new concept I’ve grasped has been akin to discovering a new star or unveiling a nebula's secrets. This journey, filled with vast opportunities for exploration and discovery, mirrors the quest for understanding the universe. It’s a journey I am excited to continue at Harvard.


From my earliest memories, the celestial canvas captivated me, its stars and galaxies sparking an insatiable desire to unravel cosmic mysteries. A decade ago, I embarked on a personal stargazing journey armed with my own telescope. Every snapshot taken—from the intricate details of Orion and M31 to the elusive NGC objects—underscored the challenges and determination inherent in astrophotography. The quest for clear skies, often leading me to the solitary expanse of the desert, only intensified my motivation to pursue my dreams in astrophysics. Each effort, each frame captured, speaks to the tenacity driving my pursuit of understanding the cosmic wonders that have captivated me.


During my undergraduate years, I captivated by the concept of dark matter, a mysterious entity accounting for most of the universe's mass yet remaining hidden from direct observation. This fascination led me to start Astronomy projects. During my Avionics studies, I used what I learned in my coursework in antenna and microwave to build my own VLF radio telescope. While this project faced challenges and was not entirely successful due to limited technical resources, it ignited a deep passion for astronomical data. This experience steered me towards observational astronomy.


While studying Space Engineering, I specialized in Astrodynamics, Statistics, and Deep Learning. My passion was ignited during astrophysics projects, including the determination of the hydrogen-alpha wavelength and calculating the redshift factor and velocity, revealing fascinating insights into stellar motion. I worked on the fuzzy-based Torino Scale, to enhance the accuracy of impact risk assessments for near-Earth objects. During my coursework in Orbital Mechanics, I delved into numerous projects, ranging from orbit determination and simulation to advanced topics including celestial mechanics and perturbation theory. I honed my analytical skills by tackling complex scenarios involving multi-body gravitational interactions, orbital maneuvers, and rendezvous strategies. These experiences motivated me to conduct research on integrating relativistic mechanics to address intricate relativistic effects, including time dilation, length contraction, and relativistic Doppler shift, particularly in the context of interstellar missions. 


My interest in Dark Matter and ISM stems from reading CFA scientists' paper on a three-dimensional map of dust reddening, using Gaia, Pan-STARRS 1 and 2MASS. This motivated me to explore AI for ISM mapping. This was the turning point at which I found the main topic I always wanted. I have found the crossroads where all my ideas and dreams converge. This is a place where I can leverage my diverse background in engineering, statistics, machine learning, and astronomy to push the boundaries of knowledge.

I aim to accelerate cosmological inference by integrating statistics and machine learning in astrophysics. My focus is on deciphering meaningful signals from vast astrophysical datasets, employing deep learning as surrogate models. In Computational Astrophysics, I'm enthusiastic about merging Astroinformatics, Astrostatistics, and AI to navigate extensive datasets from astronomical surveys. Notably, I deploy machine learning to enhance simulations of the universe's large-scale structure, advancing our understanding of the cosmic web and voids.

Fascinated by the Cosmic Microwave Background (CMB), I explore denoising diffusion probabilistic models (DDPMs) to unravel the early universe's intricacies. Leveraging innovative tools like equivariant wavelets, I aim to generate realistic astrophysical fields from cosmological simulations. In weak gravitational lensing, I plan to use Deep Neural Networks (DNNs) to exploit distortions in distant galaxies' morphologies induced by large-scale structures, enhancing our comprehension of the universe and enabling the detection and classification of gravitational waves.

In ISM mapping, I'm intrigued by the potential of diffusion models and generative adversarial networks (GANs). I intend to use GAN models to map interstellar dust emission's intensity and infer its statistical properties from observational data. Leveraging my engineering expertise, I aim to develop Digital Twins of ISM using deep learning techniques, creating a high-fidelity replica to detect subtle anomalies indicative of phenomena such as dark matter presence.

With my expertise in Large Language Models (LLMs), I aim to enhance AstroBERT’s generative capabilities and further optimize AstroLLaMA’s perplexity and improve their performance in astronomy-related tasks.


I'm drawn to Harvard not just for its cutting-edge research or renowned faculty but for its dynamic blend of innovation and tradition, creating a unique academic ethos. Beyond its global reputation, Harvard's legacy of uncovering the extraordinary in the ordinary propels my academic pursuits. The luminaries here are more than educators; they embody an insatiable quest for knowledge, serving as mentors and collaborators in a shared odyssey. The convergence of their expertise with my research interests is not just compelling; it sparks a symbiotic intellectual journey. Amidst diverse research labs, I see not just opportunities but a canvas to shape my narrative in academia. What captivates me most is the potential for unconventional camaraderie, a synergy beyond ordinary collaborative bounds. In this vibrant community, I envision not just colleagues but kindred spirits propelling me toward an uncharted trajectory, fostering a profound sense of belonging in the global Astronomy community.

As a professor, I'm dedicated to fostering an inclusive STEM community, empowering the next generation of scientists, regardless of their background. Through my passion for teaching, I aim to create an environment where everyone thrives. This commitment is not just professional; it's a personal mission. As a first-gen student, I faced academic challenges which honed my resilience, time management, and leadership skills. It inspired a passion for advocating diversity and equity, leading me to mentor students globally. I'm prepared for Harvard's challenges.



\vspace{0.05cm}
Thank you for considering my application.\vspace{-0.1cm}

Sincerely, \vspace{-0.1cm}

Arman Asgharpoor Golroudbari


\end{large}

\end{document}